\documentclass[11pt]{article}
\usepackage[margin=1in]{geometry}
\usepackage{amsmath,amssymb,array}
\usepackage[T1]{fontenc}
\setlength{\emergencystretch}{3em}
\usepackage[hidelinks]{hyperref}
\hypersetup{
  pdftitle={Supplementary Material: Formalization Scope and Explicit-Example Checks},
  pdfauthor={Myunglae Jo}
}
\title{Supplementary Material\\Formalization Scope and Explicit-Example Checks}
\author{Myunglae Jo}
\date{}
\begin{document}
\maketitle

\section*{S1. Purpose}
The companion source package verifies selected finite combinatorial components
of \emph{Small Rotation Blocks and the Order-Six Stable Marriage Problem}.
It is not an end-to-end Lean formalization of stable marriage or of the
statement $f(6)=48$.  The manuscript supplies the complete upper-bound argument
and cites the published lower-bound construction.

\section*{S2. Fixed-token combinatorics}
The module \texttt{F6SmallBlockTokenCapacity.lean} uses fixed finite types
of participants $B$ and tokens $W$, bijective states $B\simeq W$, supports
of size at least two, unchanged nonparticipants, and strictly worsening
participants.  Its principal declarations are
\begin{verbatim}
TokenRun.participation_count
TokenRun.TWO_SITE_CAPACITY
TokenRun.THREE_SITE_CAPACITY
TokenRun.SMALL_BLOCK_TOKEN_CAPACITIES
TokenRun.SIX_SITE_INCIDENCE_OPTIONAL
\end{verbatim}
in namespace \texttt{F6SmallBlockTokenCapacity}.  Writing $q$ for the number
of transitions in a run, they prove the participation bound $|B|-1$, the
capacities $q\le1$ for $|B|\le2$ and
$q\le2$ for $|B|\le3$, and $q\le15$ for $|B|\le6$.  The three-participant proof includes incidence
equality, forced pair swaps, and token return.

\section*{S3. Abstract support geometry}
The module \texttt{F6RotationSupportCI2.lean} has an abstract finite poset
$R$ and a support map to a finite set of at most six participants.  It assumes
support size at least two and comparability of elements whose supports
intersect.  The capacities for convex subsets using at most two and three
participants are the explicit hypotheses \texttt{HCAP2} and \texttt{HCAP3}.
From these hypotheses, the declaration
\begin{verbatim}
F6RotationSupport.widthLE3_and_ci2_of_rotation_support
\end{verbatim}
proves width at most three and CI2.  The source package does not construct actual
stable-marriage rotations or formally connect their execution to the abstract
token model; that connection is made in the manuscript.

\section*{S4. Finite-poset counting}
The modules \texttt{F6ConvexWindow.lean},
\texttt{F6ConvexWindowFull.lean}, and
\texttt{F6ConvexWindowN15.lean} establish the pair and triple estimates and
the fifteen-element refinement.  The final declaration is
\begin{verbatim}
F6ConvexWindow.totalAntichains_le_48
\end{verbatim}
with hypotheses \texttt{WidthLE 3}, \texttt{CI2}, and
\texttt{Fintype.card V <= 15}.  No preference-profile classification is an
input to this theorem.  The chain decomposition is supplied by the bundled Lean proof of Dilworth's theorem.

\section*{S5. Dilworth source and reproducibility}
The bundled \texttt{DilworthTheorem.lean} has source repository
\begin{quote}\small
\url{https://github.com/sneed-and-feed/lean-theorems-1}
\end{quote}
with frozen commit
\begin{quote}\ttfamily\small
4fce2c92ad4321064a06afdc86d0272bc625ce14
\end{quote}
and SHA-256
{\footnotesize
\begin{verbatim}
bfbd7c9505c0bdc4d8ba7ff49b3a1b39243ee35f3ad21d87a287a3df3a48faea
\end{verbatim}
}
The package records the CC0-1.0 dedication associated with that source
repository.

The frozen v07 \texttt{LEAN\_FORMALIZATION\_REVIEWER\_PACKAGE.zip}
has SHA-256
{\footnotesize
\begin{verbatim}
0d6524cf50fa8dfe34c456ebeb1476add48e4e39bbd00ec37938a2f4b912fa63
\end{verbatim}
}
It pins Lean 4.27.0 and mathlib v4.27.0; the mathlib commit is
\begin{quote}\ttfamily\small
a3a10db0e9d66acbebf76c5e6a135066525ac900.
\end{quote}
The public source and reproducibility repository is
\url{https://github.com/insammael/stable-marriage-f6}.
In that repository, the audited source files are exposed
directly under \texttt{lean/} and migrated to a Palomar-supported Lean
toolchain, with the toolchain, Mathlib revision, and transitive dependencies
pinned by \texttt{lean-toolchain}, \texttt{lakefile.toml}, and the committed
\texttt{lake-manifest.json}.  From the public-repository root, run
\begin{verbatim}
lake build
lake env lean palomar/AxiomAudit.lean
\end{verbatim}
The v08 release audit records the build and axiom outputs and separately
checks the migrated source against the frozen v07 source.
The reported load-bearing axioms are \texttt{propext},
\texttt{Classical.choice}, and \texttt{Quot.sound}.

\section*{S6. Components outside the formalization}
The restricted-instance observation, the explicit extremal fifteen-element
poset, the necessary equality conditions, the $DS_3$--$P_*$ isomorphism,
the general order-$N$ support-budget statement, and the
stable-matching--downset correspondence remain outside the bundled
formalization.  The upper bound does not rely on the auxiliary
example calculations below.

\section*{S7. An explicit extremal poset}
For $P_* = \{1,\ldots,5\}\times\{1,2,3\}$, put
\[
(i,a)<(j,b)\quad\Longleftrightarrow\quad
j\ge i+2\ \text{or}\ (j=i+1\text{ and }a\ne b).
\]
The manuscript directly proves that $P_*$ has width three and CI2.

Its antichain counts are $(1,15,27,5)$, totaling $48$.  An auxiliary check
enumerated all $2^{15}$ subsets, tested pairwise incomparability, and tested
CI2 for every comparable pair.  It returned the same counts, no larger
antichains, and maximum common-incomparability size two.  This independently
checks the explicit extremal example.

\section*{S8. A classical order-six lower-bound instance}
The preference instance used by the auxiliary checker is the $DS_3$
construction of Benjamin, Converse, and Krieger \cite{BenjaminCK1995}:
$DS_n$ is defined in their Section~3 (p.~287) as a $2n\times2n$ Latin
square built from $S_n$, with an algebraic definition of its entries;
Fig.~1 on that page displays $DS_3$ explicitly, and Section~5 (p.~291)
states that $DS_3$ produces $48$ stable matchings.
The table below is a transcription of Fig.~1 and the Section~3 algebraic
definition.  The publication identifies and defines the instance; the
checker verifies the mathematical properties of the transcribed input.
The row indexed by a man gives his rank of each woman, with rank $1$ best:
\[
R=\begin{pmatrix}
1&3&5&2&4&6\\
3&5&1&6&2&4\\
5&1&3&4&6&2\\
2&4&6&1&3&5\\
6&2&4&3&5&1\\
4&6&2&5&1&3
\end{pmatrix}.
\]
Woman $j$ assigns rank $7-R_{ij}$ to man $i$.  Enumerating all $6!=720$
permutations $p$ and rejecting any one for which there is a pair $(i,j)$
satisfying
\[
R_{ij}<R_{i,p(i)}
\quad\text{and}\quad
7-R_{ij}<7-R_{p^{-1}(j),j}
\]
returns $48$ stable matchings.  The bundled checker also
recovered $15$ join-irreducibles of the complete stable-matching lattice,
with antichain counts $(1,15,27,5)$, and found all $36$ pairs attainable.
All $15$ join-irreducible cover moves have two-man support, and each man
participates five times.

\subsection*{S8.1. Explicit isomorphism certificate}
Let $r_1,\ldots,r_{15}$ be the join-irreducibles in the deterministic order
of the bundled checker: stable matchings are sorted by the sum of men's
ranks, then lexicographically, and those with exactly one lower cover are
retained in that order.  Define
\[
\phi(r_{3(i-1)+k})=
\begin{cases}
(i,k),&i\text{ odd},\\
(i,4-k),&i\text{ even},
\end{cases}
\qquad 1\le i\le5,\quad 1\le k\le3.
\]
The five levels, written as indices, are
\[
\{1,2,3\},\quad\{4,5,6\},\quad\{7,8,9\},\quad
\{10,11,12\},\quad\{13,14,15\}.
\]
Covers occur only between consecutive levels.  Each such cover graph is
$K_{3,3}$ minus a perfect matching; the four missing matchings are
\[
\begin{aligned}
&\{(1,6),(2,5),(3,4)\},\\
&\{(4,9),(5,8),(6,7)\},\\
&\{(7,12),(8,11),(9,10)\},\\
&\{(10,15),(11,14),(12,13)\}.
\end{aligned}
\]
The checker verifies for every pair $u,v$ that
\[
r_u\le r_v\quad\Longleftrightarrow\quad
\phi(r_u)\le_{P_*}\phi(r_v).
\]
Writing $J(DS_3)$ for the join-irreducible poset of the stable-matching
lattice, this gives $J(DS_3)\cong P_*$.
For a finite poset $R$, the join-irreducible elements of its ideal lattice
$\mathcal J(R)$ are precisely the principal ideals $\downarrow r$, so
$r\mapsto\downarrow r$ identifies $R$ with
$\operatorname{Ji}(\mathcal J(R))$.
Combining this identification with the classical correspondence between
the stable-matching lattice and $\mathcal J(R)$, where $R$ is the rotation
poset (see the Irving--Leather reference in the manuscript), shows that
this join-irreducible poset is isomorphic to the rotation poset of $DS_3$.
This finite certificate is auxiliary to the universal upper-bound proof.

\section*{S9. Reproducing the auxiliary checks}
The supplementary package includes \texttt{AUXILIARY\_CHECKS/}, containing
\texttt{check\_ds3\_pstar.py}, its expected \texttt{DS3\_PSTAR\_CHECK.json},
a README, and SHA-256 checksums.  From that directory run
\begin{verbatim}
python3 check_ds3_pstar.py
\end{verbatim}
The checker uses only the Python standard library.  The command regenerates\newline
\texttt{DS3\_PSTAR\_CHECK.json} and must reproduce the bundled expected JSON
exactly, byte-for-byte; its expected hash is recorded in the README and
\texttt{SHA256SUMS.txt}.

\begin{thebibliography}{9}
\bibitem{BenjaminCK1995}
A. T. Benjamin, C. Converse, and H. A. Krieger,
How do I marry thee? Let me count the ways!,
\emph{Discrete Applied Mathematics} 59 (1995), 285--292.
doi:10.1016/0166-218X(95)80006-P.
\end{thebibliography}
\end{document}
